% Unnumbered local research note. No global atlas, integration, or movement.
\subsection*{Two-end formal pairs and the actual scalar coefficient gerbe}

Let $k$ be algebraically closed of characteristic $p>0$, and let
$N\ge2$.  The statements below concern the completed $A_{N-1}$
singularity and its actual minimal resolution.  They do not assert that
an arbitrary plateau, a forked RDP block, or a model with additional
$(-1)$-curves is this pair.  Every scalar coefficient order $n$ and
root order $L$ used below is prime to $p$; $N$ need not be.

\paragraph{Exactly two end branches can be straightened as formal pairs.}
Write
\[
 S=k[[x,y]],\qquad
 A=S^{\mu_N}=k[[u,v,w]]/(uv-w^N),\qquad
 u=x^N,\quad v=y^N,\quad w=xy,
                                                        \tag{C.TP.1}
\]
where $\mu_N$ has weights $(1,-1)$.  Suppose two smooth formal
branches meet the two ends of the exceptional chain transversely and
away from its nodes.  If $N=2$, their attachment points on its single
exceptional curve are required to be distinct.  There are equivariant
formal coordinates in which these two branches are the two axes.

Here is a construction retaining their entire formal expansions.
At the right end the chart is
$v=\eta$, $w=\eta\tau$, $u=\eta^{N-1}\tau^N$.
The right branch is $\tau=f(v)$ for $f\in k[[v]]$.
The triangular equivariant automorphism
\[
        x'=x-f(y^N)y^{N-1},\qquad y'=y                 \tag{C.TP.2}
\]
sends that branch to $x'=0$.  It has an explicit inverse, so no
separability or division by $N$ is involved.  For $N>2$ it leaves
the left-end attachment chart unchanged to leading order.  The
transformed left branch can therefore be written
$w'/u'=g(u')$, where $u'=(x')^N$ and $w'=x'y'$.
Now apply the second triangular automorphism
\[
        x''=x',\qquad
        y''=y'-g((x')^N)(x')^{N-1}.                   \tag{C.TP.3}
\]
It sends the left branch to $y''=0$ and preserves $x'=0$.
For $N=2$, first choose linear coordinates so the right attachment
direction has $y\ne0$.  After (C.TP.2), the distinct left direction
has $x'\ne0$, and (C.TP.3) applies.  The $\mu_2$-action is scalar
also in characteristic two, so the preliminary linear change is
equivariant.  Both triangular maps preserve the two weights because
their nonconstant factors have weight zero and $N-1\equiv-1\pmod N$.
Thus the construction is valid even when $p\mid N$.

This is stronger than choosing a parameter nonzero on a finite
attachment set: it straightens the specified two formal branches.
It says nothing about simultaneously straightening additional branches
or retaining a previously fixed coordinate presentation.

\paragraph{The bare two-axis Brauer class.}
Put $U=\operatorname{Spec}A[1/w]$.  Choose compatible identifications
of prime-to-$p$ roots of unity with constant cyclic coefficients.
Then
\[
       \operatorname{Pic}(U)=0,\qquad
       \operatorname{Br}(U)[n]\simeq\mathbb Z/n,
       \qquad H^2(U,\mu_n)\simeq\mathbb Z/n.          \tag{C.TP.4}
\]
An orientation of the two axes identifies the generator with the
symbol $(u,w)_n$; reversing conventions replaces it by its inverse.

For the Picard assertion, extend a line on $U$ to a rank-one
reflexive module on the normal surface $A$.  The already checked
$A$-classification gives the ideals $I_a=(u,w^a)$, $0<a<N$, and
$I_0=A$.  All become trivial on $U$, where $u,w$ are units.
Let $Y\to\operatorname{Spec}A$ be the minimal resolution.
Proper strict-local base change gives
$H^2(Y,\mu_n)\simeq (\mathbb Z/n)^{N-1}$ on the exceptional
tree.  The full lines $F_a$ of $I_a$ have the elementary
multidegree vectors, so their Chern classes generate this group.
Kummer therefore gives $\operatorname{Br}(Y)[n]=0$.

Purity now injects $\operatorname{Br}(U)[n]$ into the residues
on the two strict axes and the exceptional chain.  Residues on each
exceptional $\mathbb P^1$ are unramified away from its two boundary
marks.  Such a prime-to-$p$ character is determined by one cyclic
value, with opposite values at the two marks.  A strict formal axis
has function field $k((t))$ and the same one-parameter Kummer
description.  The secondary-residue relation at each node propagates
one value along the whole path.  Thus restriction to the first axis
is injective and has at most $n$ values.  The symbol $(u,w)_n$ has
there the Kummer residue of the uniformizer $u$, of exact order $n$.
It realizes all values.  This proves the Brauer assertion; Picard
vanishing and Kummer prove the last assertion.  Only injectivity and
the secondary-residue relation are used, not a claim of exactness
at every term of an unproved Gersten complex.

The purity input is the previously used absolute-purity theorem of
Riou in the ILO volume, with its divisor and normal-crossings
consequences \cite[Exp.~XVI, \S3.1]{A_ILO2014}.  Precise fresh and
reused source scopes appear at the end of this note.

\paragraph{The punctured root stack retains more than the bare class.}
For a common root order $L$ define the diagonalizable group $G_L$ by
\[
 X^*(G_L)=\mathbb Z^2/
       \langle(NL,0),(0,NL),(L,L)\rangle.
                                                        \tag{C.TP.5}
\]
It acts on $S$.  Its quotient is $A$, now with
$u=x^{NL}$, $v=y^{NL}$, $w=(xy)^L$; its subgroup $\mu_N$
still acts with weights $(1,-1)$.  Put
\[
 T=S^{\mu_N},\qquad Q=G_L/\mu_N\simeq\mu_L^2,
 \qquad \mathcal U_L=[(\operatorname{Spec}S\setminus\{0\})/G_L].
                                                        \tag{C.TP.6}
\]
The last stack is the order-$L$ root stack along the two axes on
the punctured singularity.  Let $\alpha$ be an \emph{actual}
$\mu_n$-gerbe on this stack.  Its restriction to the bare $U$ has
some Brauer order $h$.  Extension to $\mathcal U_L$ forces
$h\mid L$, and of course $h\mid n$.

There is a finite central extension
\[
       1\longrightarrow\mu_n\longrightarrow H
                    \longrightarrow G_L\longrightarrow1       \tag{C.TP.7}
\]
whose quotient gerbe on $S\setminus\{0\}$ is $\alpha$.
To see the actual descent, purity in codimension two on the strictly
local regular surface gives
$H^1(S\setminus\{0\},\mu_n)=H^2(S\setminus\{0\},\mu_n)=0$.
Choose a neutral object and trivialize its Isom torsor on
$G_L\times(S\setminus\{0\})$.  The same vanishings hold on
every nerve term: $G_L^j$ is a finite disjoint union of local
Artinian $k$-schemes, and prime-to-$p$ etale cohomology is unchanged
by their nilpotent thickenings.  The remaining composition cocycle
is a $\mu_n$-valued function on $G_L^2\times(S\setminus\{0\})$.
Such a function is constant on each component in the second factor,
so is a function on $G_L^2$.  Its cocycle identity defines (C.TP.7).
This argument keeps the group scheme when $p\mid N$; it does not
replace $G_L$ by its reduced subgroup in the representation category.

Equivariant reflexive extension and equivariant Nakayama also give
\[
                  \operatorname{Pic}(\mathcal U_L)=X^*(G_L).
                                                        \tag{C.TP.8}
\]
For example, the extension of a line to $S$ is free, and its fibre
at the origin is a character; exact invariants lift that character
equivariantly to a basis.  Thus the difference between two
coefficient gerbes with the same bare Brauer class can retain a
nontrivial Kummer character line on $\mathcal U_L$.

\paragraph{The scalar band does not enlarge the list of simple types.}
Write a diagonal element of $G_L$ as $(\zeta,t)$, acting by
$(x,y)\mapsto(\zeta x,t\zeta^{-1}y)$, where
$\zeta\in\mu_{NL}$ and $t\in\mu_L$.  An alternating
$\mu_n$-valued commutator pairing is determined by one value of
order dividing $\gcd(L,n)$.  It kills $\mu_N$ and factors through
the displayed $Q$.  For commutator order $h$ its radical is
\[
 K=\ker\bigl(G_L\longrightarrow\mu_h^2\bigr),
 \qquad |K|=N(L/h)^2.                                  \tag{C.TP.9}
\]
Here the map is the composite $G_L\to Q\to\mu_h^2$.
The primitive Heisenberg extension of $\mu_h^2$, pushed out to
$\mu_n$, has bare symbol order $h$ by (C.TP.4).  Subtracting it
from (C.TP.7) makes the commutator trivial.  A commutative extension
of the indicated diagonalizable group by $\mu_n$ is itself
diagonalizable: its prime-to-$p$ part is finite etale abelian and
its $p$-primary part is unchanged by the prime-to-$p$ kernel.
Dualizing its character sequence identifies it with the root of
a character line.  Such a line is trivial on $U$ by (C.TP.4).
Consequently the commutator order really is the bare Brauer order
$h$, and the residual discrepancy is a Kummer line, not an extra
unseen Brauer period.

The preimage $H_K$ of $K$ is central and diagonalizable.  The
characters of $H_K$ restricting to scalar weight one on $\mu_n$
form a nonempty torsor under $X^*(K)$, so there are exactly $|K|$
of them.  For each such character, the residual twisted group
algebra of $H/H_K\simeq\mu_h^2$ is $\operatorname{Mat}_h(k)$:
its commutator is nondegenerate, and the shift-clock matrices give
all $h^2$ matrix units.  Thus there is one simple of dimension $h$
for each character.  In particular,
\[
 \#\{\text{scalar-weight-one simple types}\}=N(L/h)^2,
 \qquad \dim(\text{each simple})=h.                    \tag{C.TP.10}
\]
This number is independent of $n$.  A character of $H_K$ need
\emph{not} extend to all of $H$.  No such extension, splitting of
the scalar inertia, or independent choice of phase linearizations
is used here.

The same calculation classifies actual geometric punctured bundles.
Their pullbacks to $S\setminus\{0\}$ extend reflexively to free
$S$-modules.  Hartogs extends the action and its inverse after
tensoring with $k[H]$, since $x,y$ remain a regular sequence there;
the group identities extend after the second such tensor product.
The group $H$ is linearly reductive, because invariants for its
normal $\mu_n$ and diagonalizable quotient $G_L$ are exact.
An equivariant section of the fibre at the origin now gives the
actual form $S\otimes V$.  Multiplicities of the types (C.TP.10)
are therefore intrinsic.  No relative reflexive hull over a
nonreduced coefficient base is asserted in this step.

\paragraph{Neutralization on the $\mu_N$ cover and its remaining coherence.}
Let $H_N$ be the preimage of $\mu_N$.  It is central and
diagonalizable.  Choose a character $\lambda_0$ of $H_N$ of
scalar weight one.  On the gerbe above
$T^*=\operatorname{Spec}T\setminus\{0\}$ this gives an actual
neutralizing twisted line $\Lambda$.  Its $n$th power descends
to the character line $I_c|_{T^*}$ for some $0\le c<N$:
\[
                      \Lambda^n\simeq I_c|_{T^*}.     \tag{C.TP.11}
\]
The selected isomorphism is part of the data.  Tensoring by
$\Lambda^{-1}$ is performed only on this cover.  It is not a
claim that $\lambda_0$ extends to $H$.

There is a pointwise full extension on a chosen coefficient gerbe
over the toric resolution $Y_T$.  Extend $I_c$ by its ordinary
full line $F_c$ and take $\sqrt[n]{F_c/Y_T}$, with its tautological
line $\widehat\Lambda$.  The neutralized module of a punctured
object is a direct sum of the $I_a$, so extend it by the corresponding
sum of $F_a$ and then tensor by $\widehat\Lambda$.

The quotient $Q$-coherence must be retained in doing this.
Choose lifts $t_g\in H(k)$ for the finite etale group $Q$.
Their multiplication defects
$c(g,g')=t_gt_{g'}t_{gg'}^{-1}\in H_N(k)$ are central.
On $T^*$ the chosen line $\Lambda$ has the resulting weak
$Q$-transport; its scalar comparison factors are
$\lambda_0(c(g,g'))$, with the inverse if pullback conventions
are reversed.  Their $n$th powers are the corresponding transport
factors of $I_c$.  The latter line maps and their inverses extend
to $F_c$ by full Hom, because the diagonal toric deck maps preserve
the reflexive character modules.  Use these extended line maps
and the \emph{same} scalar comparison factors to define the weak
action on $\sqrt[n]{F_c/Y_T}$.  Its associativity identities hold
because the original central-extension identities hold, and equalities
of line maps extend from the puncture.

For each neutralized object, the actual deck maps between its
ordinary full summands extend by full Hom as well.  The inverse
scalar defects cancel those of $\widehat\Lambda$, exactly as
on the puncture.  This gives a full twisted object on the quotient
by this weak action, with its original punctured descent data.
Maps extend by the same argument.  For two objects of the same
scalar weight, the common $\widehat\Lambda$ cancels in Hom;
ordinary full-Hom vanishing and exact $Q$-invariants apply.
Actual Artinian bundle families on this fixed chosen extension
therefore have the previously proved full-object rigidity.  This
last assertion is about actual bundle families, not flatness of a
relative fullization construction.

This constructs a particular extension of the given punctured
coefficient gerbe.  It is \emph{not} an identification with every
previously chosen gerbe extension on the full resolution.  In
particular its comparison to that extension, or its descent to
specified unequal end-root orders, must retain the actual Kummer
line and redundant-inertia data.  The object-level end-descent
test in the separate B mixed-end note applies only after that
coefficient-gerbe comparison is supplied.

\paragraph{The coefficient line has bounded one-step phase corrections.}
Suppose $n=p^f-1>1$, and retain the actual power identification
in (C.TP.11).  Set
\[
 s_a\equiv p^a\pmod n,\quad 0\le s_a<n,\qquad
 \Lambda_a=\Lambda^{p^a}\otimes
                  F_c^{-\lfloor p^a/n\rfloor}.
\]
All lines here are interpreted on the chosen root gerbe, and the
same weak deck data are transported with these tensor operations.
Then the actual identities are
\[
 \Lambda_a^p\simeq\Lambda_{a+1}\otimes F_c^{k_a},
 \qquad k_a=\lfloor ps_a/n\rfloor\in\{0,\ldots,p-1\},
 \qquad \Lambda_f\simeq\Lambda_0.                    \tag{C.TP.12}
\]
The last identity uses $\Lambda^n=F_c$; it is not a new arbitrary
cyclic linearization.  For $n=1$ the band is trivial and this residue
notation is omitted.

For an ordinary summand $F_i$, a one-step comparison therefore
involves $F_i^pF_c^{k_a}$.  Write $pi+k_ac=qN+b$, $0\le b<N$.
Its actual reflexive target is the fractional module $U^qI_b$,
where $U=x^N$.  At the $j$th exceptional curve the effective
fullization discrepancy is
\[
 d_j=p\min(j,i)+k_a\min(j,c)-qj-\min(j,b),
 \qquad 0\le d_j\le(2p-1)\min(j,N-j).                \tag{C.TP.13}
\]
The lower bound is the inclusion of the product ideal in its
reflexive hull.  For the upper bound, the piecewise-linear function
vanishes at $0,N$ and every slope has absolute value at most
$p+k_a\le2p-1$.  The scalar $U^q$ and all actual representation
intertwiners are retained in this comparison.

When the contracted block really lies on the controlled model of
the existing Jacobian calculation, that calculation gives at every
vertex $\min(j,N-j)\le R_j\le\lambda C_j$.
Thus (C.TP.13) has a fixed inherited-$C$ pole bound.  This invokes
the actual Jacobian hypothesis; it is false to replace $C$ by the
reduced fundamental cycle uniformly in $N$.  Neither the scalar
band size $n$ nor an expanded $f$-fold word enters the one-step
estimate.

\paragraph{Remaining identification and quantitative gaps.}
The preceding statements concern one exact completed two-end
$A_{N-1}$ pair, one fixed actual punctured coefficient gerbe, and
a chosen full extension with its retained weak deck data.  They
do not supply a global, all-coefficient replacement functor.
An application still needs compatibility with the neighboring core,
the previously fixed full gerbe and order action, and unequal
end-root structures.  It must also treat all large normal-thickness
overlaps.  The first coarse end-jet bound of the separate fat-end
note cannot simply be repeated at every thickness: even the
ordinary structure sheaf has a quadratic thick-end defect.
A small actual object presentation on the core and a coupled
whole-cycle gluing estimate remain open.  Forked RDP blocks and
additional blowup components have not been identified with this
two-end model.  No unbounded-period hereditary count or strict
main-gate reduction follows from this note.

\paragraph{Exact source and proof scope.}
The sequential formal changes, group schemes, central-extension
refinement, representation count and scalar phase formulas are
written out here.  The complete B common-root and mixed-end quantum
notes and the complete existing full-Hom and Artinian-framing proofs
were freshly read.  The regular-surface reflexive-extension and
equivariant grading inputs are reused in their already audited scopes.
For purity, the official Riou draft of 14 November 2016,
\url{https://www.cmls.polytechnique.fr/perso/orgogozo/travaux_de_Gabber/16-classes-de-Chern/gysin.pdf},
was freshly opened: the exact Theorem 3.1.1 statement and the
complete written proofs of Corollaries 3.1.3 and 3.1.4 on printed
pages 15--16 were read.  The full proof of Theorem 3.1.1 and its
recursive inputs were not freshly reread in this interval; they
remain the known previously used absolute-purity input.  The
normal-crossings corollary is 3.1.4 in this draft, whereas its
number differs in the retained older source.  No new external
classification of ramified orders is invoked.  This is an unnumbered,
unpromoted research note.
